\chapter{Investigación de Operaciones: repartir el presupuesto}
\label{cap:io}
\textit{Criterio 6 de la rúbrica (12\,\%): formula variables de decisión, función objetivo y restricciones; la solución es
justificable.}

\section{La pregunta}
¿Cuánto dinero de la campaña va a cada lugar del sur, en qué mes y por qué canal, para traer el mayor número de visitantes
\textbf{sin anunciar donde ya está lleno} y \textbf{sin rebasar la capacidad} de nadie? Responde al incidente de
investigación de operaciones: qué es ``óptimo'', qué priorizar si los objetivos chocan, qué restricciones son
indispensables y qué pasa si se cambian.

\section{Los parámetros: cuánto rinde un peso en cada canal}
\subsection{El embudo de un anuncio}
\begin{intuicion}
Un anuncio funciona como un embudo: muchas personas lo ven (impresiones), algunas le dan clic, y de esas, algunas planean
el viaje (conversión). Se paga por clic. Lo que importa es cuántas conversiones trae cada peso.
\end{intuicion}

\begin{gather*}
\text{clics}=\frac{\text{pesos}}{\text{costo por clic}},\qquad \text{conversiones}=\text{clics}\times\text{tasa de conversión},\\
r_c=\frac{\text{conversiones}}{\text{peso}}=\frac{\text{conversión}_c}{\text{costo por clic}_c\ (\text{dólares})\times\text{pesos por dólar}} .
\end{gather*}

\textbf{Datos} (categoría \emph{Travel}, WordStream 2025; promedios de anunciantes de EE.\,UU., usados como supuesto
etiquetado): Google, clic de \$2.12 dólares y conversión de 5.75\,\%; Facebook, clic de \$0.51. El dólar es el promedio de
agosto de 2026 en FRED: 17.0609 pesos (21 días con dato).

\begin{amano}[ · conversiones por peso]
Google: $r=\dfrac{0.0575}{2.12\times17.0609}=\dfrac{0.0575}{36.169}=0.001590$, es decir, \textbf{1.59 visitantes por cada
\$1,000}.

Facebook: $r=\dfrac{0.0575}{0.51\times17.0609}=\dfrac{0.0575}{8.701}=0.006608$, es decir, \textbf{6.61 por cada \$1,000}.

Facebook rinde cuatro veces más por peso porque su clic es cuatro veces más barato (con la misma conversión supuesta).
\end{amano}

\textbf{Hueco declarado}: la conversión de Facebook para turismo \emph{no está publicada}. El caso base usa la de Google
Travel (5.75\,\%) y la sensibilidad prueba 3\,\% y 6.38\,\% (la mediana de Facebook en todas las industrias).
\textbf{Supuesto}: cada conversión cuenta como un visitante. Es una cota alta: si llegan menos, la capacidad queda aún más
protegida.

\section{El modelo completo}
\subsection{Conjuntos y variables}
\begin{itemize}
\item Meses $m$: de octubre de 2026 a junio de 2027, los 9 que tienen calendario y escenarios para los tres lugares.
  Después no hay pronóstico, y no se extrapola.
\item Lugares $l$: Chetumal, Bahía, Ruta. Canales $c$: Google, Facebook.
\item Escenarios $s$: malo, probable y bueno (percentiles 10, 50 y 90 del Monte Carlo), con pesos $p_s=0.3,\ 0.4,\ 0.3$
  (regla de Swanson para resumir una distribución con tres percentiles).
\end{itemize}

\textbf{Variables de decisión}:
\begin{itemize}
\item \textbf{Etapa 1} (se decide hoy, ``aquí y ahora''): $x_{m,l,c}\ge0$, pesos en el mes $m$, lugar $l$, canal $c$.
\item \textbf{Etapa 2} (recurso, se decide cuando se ve el escenario): $y_{s,m,l}\in\{0,1\}$, si el anuncio sigue
  encendido; $w_{s,m,l}\ge0$, los visitantes que sí llegan.
\end{itemize}
Visitantes que trae la campaña: $v_{m,l}=\sum_c r_c\,x_{m,l,c}$.

\subsection{Función objetivo}
\[
\max\ \sum_{s}p_s\sum_{m,l}w_{s,m,l}
\]
Los visitantes \textbf{esperados}: el promedio, ponderado por la probabilidad de cada escenario, de los visitantes que sí
llegan.

\subsection{Restricciones}
\begin{align}
&\textstyle\sum_{m,l,c}x_{m,l,c}\le B=250{,}000\times\tfrac{9}{12}=187{,}500 &&\text{presupuesto (prorrateado)}\\
&x_{m,l,c}=0\quad\text{si }(m,l)\text{ es temporada alta} &&\text{cero anuncio en temporada alta}\\
&D^{p50}_{m,l}+v_{m,l}\le C_l &&\text{capacidad probada}\\
&\textstyle\sum_{m,c}x_{m,l,c}\ge0.15\,B\quad\forall l &&\text{piso de equidad}\\
&\textstyle\sum_{m,l}x_{m,l,c}\le0.70\,B\quad\forall c &&\text{tope por canal}\\
&D_{s,m,l}+v_{m,l}\le C_l+M_{s,m,l}\,(1-y_{s,m,l}) &&\text{recurso: pausa si se rebasa}\\
&w_{s,m,l}\le v_{m,l},\qquad w_{s,m,l}\le U\,y_{s,m,l} &&\text{solo cuenta lo que llega}
\end{align}
con $U=\max_c r_c\cdot B$ (lo más que podría traer una celda) y $M_{s,m,l}=\max(D_{s,m,l}+U-C_l,\,0)$.

\subsection{Qué significa cada pieza}
\textbf{La ``M grande''.} La restricción de recurso tiene dos modos. Si el anuncio sigue ($y=1$), el término
$M(1-y)$ vale 0 y la capacidad debe respetarse en ese escenario. Si se pausa ($y=0$), el término suma $M$, que es justo lo
necesario para que la restricción ya no obligue; pero entonces $w\le U\cdot0=0$: esos visitantes se pierden. El valor de
$M$ es el más chico que logra eso, porque una $M$ demasiado grande vuelve lento al solucionador.

\textbf{Por qué ``estocástico de dos etapas''.} La decisión de hoy ($x$) se toma sin saber qué escenario va a ocurrir; la
reacción ($y$) depende del escenario. Al maximizar el promedio sobre escenarios, el modelo aprende, desde hoy, a no
gastar en meses donde un año bueno llenaría el lugar: ese dinero se perdería con probabilidad $p_{\text{bueno}}=0.3$.
Con tres escenarios, el problema se escribe como un solo modelo grande (el \emph{equivalente determinista}).

\section{Cómo lo resuelve la computadora}
\textbf{PuLP} escribe el modelo y \textbf{CBC} lo resuelve. Tiene 54 variables $x$, 81 binarias $y$ y 81 variables $w$.
\begin{itemize}
\item \textbf{Programación lineal y simplex.} Si todas las variables fueran continuas, las restricciones lineales formarían
  un poliedro convexo, y el óptimo de un objetivo lineal siempre está en una esquina. El método simplex salta de esquina
  en esquina, mejorando cada vez, hasta que ninguna vecina mejora.
\item \textbf{Ramificación y acotamiento.} Como $y$ solo puede valer 0 o 1, CBC resuelve primero la versión con $y$ continua
  (una cota optimista). Si alguna $y$ sale fraccionaria, crea dos subproblemas ($y=0$ y $y=1$) y descarta las ramas cuya
  cota no puede superar la mejor solución entera encontrada.
\item Tarda \textbf{0.3 segundos}.
\end{itemize}

\subsection{Desempate en dos pasos (optimización lexicográfica)}
Como todos los lugares convierten igual (no hay dato de respuesta por lugar), hay muchísimas soluciones con el mismo máximo
de visitantes. El solucionador entregaba una esquina cualquiera: todo el dinero de cada lugar en un solo mes (los
\$131,250 de la Ruta, en junio).
\begin{decision}
Desempatar \textbf{proporcional al espacio libre}. Paso 1: el máximo de visitantes $z^*$. Paso 2: exigir al menos $z^*$ y,
entre esas soluciones, quedarse con la que reparte cada canal en proporción al espacio libre de cada mes:
\[
\min\sum_{m,l,c}\left|x_{m,l,c}-\pi_{m,l}\sum_{m',l'}x_{m',l',c}\right|,\qquad
\pi_{m,l}=\frac{\text{libre}_{m,l}}{\sum_{\text{permitidas}}\text{libre}},\qquad \text{libre}=1-\frac{D^{p50}}{C} .
\]
Descartados: un tope de 25\,\% por mes (una regla nueva) y dejarlo concentrado.
\end{decision}
\textbf{Cómo se vuelve lineal el valor absoluto}: se agrega una variable $d\ge0$ con $d\ge x-\text{meta}$ y
$d\ge\text{meta}-x$, y se minimiza $d$. En el óptimo, $d$ es exactamente $|x-\text{meta}|$.

\section{El resultado, a mano}
\begin{amano}[ · el máximo de visitantes]
Facebook rinde más por peso, así que se lleva su tope (70\,\%); Google, el resto:
\[
z^*=0.70\times187{,}500\times0.006608+0.30\times187{,}500\times0.001590=867.3+89.4=956.8\ \text{visitantes}.
\]
En el embudo: \$131,250 en Facebook compran $131{,}250/8.70=15{,}084$ clics y $15{,}084\times0.0575=867$ visitantes;
\$56,250 en Google compran $1{,}555$ clics y 89 visitantes. Unos \textbf{\$196 por visitante}.
\end{amano}

\begin{amano}[ · el reparto proporcional, Ruta en octubre de 2026]
Libre $=1-3{,}688/10{,}465=0.6476$. Hay 20 combinaciones de mes y lugar permitidas (las 7 restantes son temporada alta) y
sus espacios libres suman 8.0004: $\pi=0.6476/8.0004=0.0809$.
Facebook: $0.0809\times131{,}250=10{,}624$ pesos. Google: $0.0809\times56{,}250=4{,}553$ pesos. El código da lo mismo.
\end{amano}

\textbf{Reparto final}: Ruta 41.5\,\% (\$77,790), Bahía 36.9\,\% (\$69,175), Chetumal 21.6\,\% (\$40,535). Diciembre sin
anuncio en los tres lugares (temporada alta).

\begin{figure}[h]\centering
\includegraphics[width=0.85\linewidth]{f13_reparto_presupuesto.png}
\caption{Reparto del presupuesto por mes y lugar.}
\end{figure}

\section{¿Cuánto cuesta cada regla? Dualidad y precios sombra}
\subsection{Costo de cada regla}
Se resuelve el modelo sin la regla y se compara:
\begin{center}
\begin{tabular}{lrr}
\encabezado \h{Regla} & \h{Visitantes que cuesta} & \h{\%} \\
Cero anuncio en temporada alta & 0 & 0\,\% \\
Tope de capacidad probada & 0 & 0\,\% \\
Piso de equidad de 15\,\% & 0 & 0\,\% \\
Tope de 70\,\% por canal & 282 & 29.5\,\% \\
\end{tabular}
\end{center}

\subsection{Precios sombra}
\textbf{Qué son.} Todo programa lineal tiene un problema ``dual''. Su solución da, para cada restricción, un
\textbf{precio sombra} $\lambda_i=\partial z^*/\partial b_i$: cuánto mejoraría el objetivo si la restricción se aflojara
una unidad. Se calculan con las $y$ fijas en su solución (para que el modelo sea lineal).

\textbf{Holgura complementaria}: si una restricción no está ``apretada'' (le sobra holgura), su precio sombra es 0. Por eso
los precios sombra explican la tabla de arriba:
\begin{itemize}
\item \textbf{Presupuesto}: 0.00159 visitantes por peso (1.59 por cada \$1,000). Un peso extra iría a Google, porque
  Facebook ya está en su tope.
\item \textbf{Tope de Facebook}: $0.006608-0.001590=0.005019$ por peso. Cada peso que pasa de Google a Facebook suma 0.005
  visitantes. Quitar el tope suma $0.30\times187{,}500\times0.005019=282.3$, exactamente el costo de la tabla.
\item \textbf{Equidad}: precio sombra 0. Aunque en la solución del paso 1 el piso de Chetumal queda justo en 15\,\%,
  aflojarlo no trae ningún visitante: todos los lugares convierten igual, así que mover un peso de Chetumal a otro lugar no
  cambia el total. (Es un caso \emph{degenerado}: la restricción está apretada pero no estorba.) En el reparto final, cada
  lugar recibe más del 15\,\%.
\item \textbf{Capacidad}: no estorba. La campaña es chica frente al espacio que hay, y resolver sin esa regla da los mismos
  957 visitantes.
\end{itemize}

\subsection{Frontera de Pareto: visitantes contra espacio libre}
\textbf{Qué es.} Cuando hay dos objetivos que chocan (más visitantes contra menos presión), no hay un solo óptimo: hay una
\emph{frontera} de soluciones donde no se puede mejorar uno sin empeorar el otro. Se traza con el método de la
\textbf{$\varepsilon$-restricción}: se pide que ningún mes con anuncio pase de $\varepsilon$ veces su capacidad en el
escenario probable ($D^{p50}+v\le\varepsilon C$) y se maximizan los visitantes, para varios valores de $\varepsilon$.

\begin{center}
\begin{tabular}{lrrrrrrrrrr}
\encabezado \h{$\varepsilon$} & \h{100\,\%} & \h{90} & \h{80} & \h{70} & \h{60} & \h{50} & \h{45} & \h{40} & \h{35} & \h{30} \\
Visitantes & 957 & 957 & 957 & 957 & 957 & 957 & 957 & 957 & 537 & 0 \\
\end{tabular}
\end{center}
Hasta 40\,\% no se pierde ningún visitante: la campaña puede prometer que ningún mes con anuncio pasará del 40\,\% de su
capacidad probada.

\begin{figure}[h]\centering
\includegraphics[width=0.75\linewidth]{f14_frontera_pareto.png}
\caption{Frontera de Pareto.}
\end{figure}

\section{Sensibilidad: ¿y si los supuestos están mal?}
\begin{center}
\begin{tabular}{lr}
\encabezado \h{Si cambia…} & \h{Visitantes} \\
Nada (caso base) & 957 \\
Conversión de Facebook de 3\,\% & 542 \\
Conversión de Facebook de 6.38\,\% (mediana de industrias) & 1,052 \\
Clic de Facebook a \$0.42 dólares (LocaliQ 2026) & 1,143 \\
Golpe de tormenta de 25\,\% o 50\,\% & 957 \\
Presupuesto de \$500,000 al año & 1,914 \\
Presupuesto de \$1,000,000 al año & 3,827 \\
\end{tabular}
\end{center}
Los visitantes son proporcionales al presupuesto porque el modelo es lineal y la capacidad no se alcanza; el reparto entre
lugares y canales casi no cambia en ningún caso. La recomendación es estable.

\section{La etapa 2 en la práctica: el motor de la Torre}
Cada semana, para cada lugar, el motor (\codigo{envivo/motor.py}) decide: temporada alta (no se gasta), fuera del plan,
\textbf{pausado} (tormenta cerca, clima raro o llegó más gente de la esperada) o \textbf{encendido}. El dinero pausado pasa
a la siguiente semana permitida:
\begin{gather*}
b=\frac{\text{pesos del mes}}{\#\text{lunes del mes}},\qquad g_w=(b+A_{w-1})\,\ind[\text{encendido}],\\
A_w=\begin{cases}A_{w-1}+b & \text{si se pausa}\\ 0 & \text{si se enciende (se gasta todo)}\\ A_{w-1} & \text{en temporada alta o fuera del plan}\end{cases}
\end{gather*}
\begin{amano}[ · el arrastre del dinero]
Octubre tiene 4 lunes; con \$4,000 en el mes, $b=\$1{,}000$. Semana 1 pausada: $A=1{,}000$. Semana 2 encendida: se gastan
$1{,}000+1{,}000=\$2{,}000$ y $A=0$.
\end{amano}
En las 239 semanas reales: 48 pausas de lugar-semana, todo el dinero pausado se gastó después, y ``¿Ibas al norte?'' se
encendió 6 semanas.

\begin{codigoen}\codigo{campana/presupuesto.py}: \codigo{benchmarks}, \codigo{canales}, \codigo{resolver},
\codigo{precios_sombra}, \codigo{costo_de_reglas}, \codigo{pareto}, \codigo{sensibilidad}. Notebook
\codigo{05_optimizacion}.\end{codigoen}

\begin{teprofe}
\emph{``¿Qué es lo estocástico de tu modelo de optimización?''} La segunda etapa: la decisión de pausar depende de cuál de
los tres escenarios ocurra, cada uno con su probabilidad, y el objetivo es el valor esperado. \emph{``¿Cuánto cuesta la
regla ambiental?''} A este presupuesto, nada: su precio sombra es 0 porque tiene holgura.
\end{teprofe}
